\hypersetup{
  pdftitle={Kenya Multi-Hazard Catastrophe Intelligence},
  pdfsubject={Crowd intelligence, AI, catastrophe modelling, insurance, investment and resilience finance in Kenya},
  pdfkeywords={Kenya, catastrophe, flood, drought, wildfire, locust, landslide, heat, AI, actuarial, insurance, resilience finance}
}

\frontmatter
\chapter*{Series orientation}
\addcontentsline{toc}{chapter}{Series orientation}

This six-part research series defines how Kenya can combine authoritative observations, earth observation, environmental sensors, crowd intelligence, artificial intelligence and historical loss evidence to estimate multi-hazard catastrophe states and financial consequences. It covers flood, drought, wildfire and rangeland fire, locust and catastrophic crop pests, rainfall-induced landslide, severe storm and extreme heat.

The series preserves a strict boundary between evidence and authority. Model outputs are not official warnings, insurance decisions, public appropriations, accounting entries, investment recommendations or financing triggers unless a separately authorised institution and instrument make them so.

\begin{center}
\vspace{6mm}
\colorbox{SeriesPale}{\parbox{0.88\textwidth}{\centering\sffamily\color{SeriesNavy}
\textbf{Research boundary}\par
The California workshop notes contribute hypotheses, not Kenyan evidence. The actuarial charter and the Mwendo Pamoja benchmark govern method, breadth and presentation only; their subject matter is not imported.}}
\end{center}

\vspace{6mm}
\chapter*{How to use this volume}
\addcontentsline{toc}{chapter}{How to use this volume}

Part 1 states the thesis and causal architecture. Part 2 defines admissible observations. Part 3 develops spatiotemporal hazard-state models. Part 4 translates hazard, exposure and vulnerability into loss. Part 5 maps evidence into insurance, investment and resilience finance. Part 6 defines governance, validation and staged implementation. Each part remains independently citable; the combined volume provides the shared narrative and contents index.
